%%%%%%%%%%%%%%%%% Оформление ГОСТА%%%%%%%%%%%%%%%%%

% Все параметры указаны в ГОСТЕ на 2021, а именно:

% Шрифт для курсовой Times New Roman, размер – 14 пт.
\setdefaultlanguage[spelling=modern]{russian}
\setotherlanguage{english}

\setmainfont[
  Ligatures=TeX
]{Times New Roman}

\newfontfamily\cyrillicfont[
  Script=Cyrillic,
  Ligatures=TeX
]{Times New Roman}

\newfontfamily\englishfont[
  Script=Latin,
  Ligatures=TeX
]{Times New Roman}

% monofont лучше не Times New Roman
\setmonofont{Courier New}



% шрифт для URL-ссылок
\urlstyle{same}

% Междустрочный интервал должен быть равен 1.5 сантиметра.
\linespread{1.5} % междустрочный интервал


% Каждая новая строка должна начинаться с отступа равного 1.25 сантиметра.
\setlength{\parindent}{1.25cm} % отступ для абзаца


% Текст, который является основным содержанием, должен быть выровнен по ширине по умолчанию включен из-за типа документа в main.tex


%%%%%%%%%%%%%%%%%% Дополнения %%%%%%%%%%%%%%%%%%%%%%%%%%%%%%%%%

% Путь до папки с изображениями
\graphicspath{ {./Images/} }

% Внесение titlepage в учёт счётчика страниц
\makeatletter
\renewenvironment{titlepage} {
  \thispagestyle{empty}
}


% Цвет гиперссылок и цитирования
\usepackage{hyperref}
\hypersetup{
  colorlinks=true,
  linkcolor=black,
  filecolor=blue,
  citecolor = black,
  urlcolor=blue,
}


% Нумерация рисунков
\counterwithin{figure}{section}
% для сквозной нумерации рисунков
% \counterwithout{figure}{section}

% Нумерация таблиц
\counterwithin{table}{section}
% для сквозной нумерации таблиц
% \counterwithout{table}{section}

% Подписи рисунков:
% по центру, ниже рисунка, формат: Рисунок X -- Название
\captionsetup[figure]{
  name=Рисунок,
  labelsep=endash,
  justification=centering,
  singlelinecheck=true,
  position=bottom
}

% Подписи таблиц:
% справа, выше таблицы, формат: Таблица X. Название
\captionsetup[table]{
  name=Таблица,
  labelsep=period,
  justification=raggedright,
  singlelinecheck=false,
  position=top
}



% шрифт для листингов с лигатурами
\setmonofont{FiraCode-Regular.otf}[
  SizeFeatures={Size=10},
  Path = Settings/,
  Contextuals=Alternate
]

% Перенос текста при переполнении
\emergencystretch=25pt


% настройка подсветки кода и окружения для листингов
%\usemintedstyle{colorful} % делает подсветку для кода
\newenvironment{code}{\captionsetup{type=listing}}{}


% Посмотреть ещё стили можно тут https://www.overleaf.com/learn/latex/Code_Highlighting_with_minted